\section{Limitations}
\label{app:limitations}

The theory identifies when a distinction is necessary or valuable at the population-information level. It does not solve the search problem of finding the best distinction in a rich hypothesis class. The online algorithm is greedy and local: its refinement rule concerns local concepts and their nearby conflicts, not global optimality over all information structures.

The risk identities assume correctly estimated conditional distributions or conditional means. With finite data, refinement can increase estimation error, so useful population distinctions may be statistically unsafe. Complexity penalties, thresholds, validation, Bayesian priors, or minimum leaf sizes are needed in practical systems.

The finite-sample validation theorem assumes a genuinely held-out scoring buffer: fresh samples, or samples not previously used to choose proposals or update the state. Reusing the same validation buffer after previous accept/reject decisions makes the candidate class adaptive to that buffer and requires reusable-holdout or other adaptive-validation machinery beyond the Hoeffding argument stated here.

The base-information convention is restricted. If the full hidden state is made freely conditionable in the Bayes-risk benchmark, then all CA edge predicates and anchor-order events are already measurable and the sigma-algebra refinement gap vanishes. The strict partition theory therefore applies to the current runtime information or to restricted hard-route prediction classes; h-dependent soft memories are implementation approximations or restricted-class extensions.

The graph-growing CA extension specifies one canonical refinement mechanism, based on concepts and local distinction edges, but this is not claimed to be optimal. Rich domains may require learned concept-creation rules, approximate nearest-neighbor indices, nonlinear distinction edges, and better merge, forgetting, or retention-decay policies than the default budget notes in Appendix~\ref{app:dgm-extensions}. The bridge to Projection Memory and CPM is exact for one-hot finite partition memories. Continuous neural representations need additional modeling choices: when a continuous conflict should create a new concept, how local projection memories should be inherited, how retention time scales should be allocated, and how refined information structures interact with differentiable backbones.

CA's dense read avoids sparse candidate-recall assumptions at the concept budgets used in the main experiments, and the current implementation includes fused kernels for the dense concept operations and chunk scan. Large-budget CA will still need sparse concept retrieval, sharded concept banks, and additional fused write/read kernels; at that point candidate recall again becomes an engineering assumption that should be measured in large systems. Edge gates can reject incompatible candidates in the graph-growing extension, but they cannot select a compatible concept that the retrieval layer never exposes.

The large-model layer is an architectural path rather than a completed scaling theorem. It gives a trainable read path, a continuous write path, a trainable retention decay, and a default budget policy with explicit complexity, but it does not prove optimization guarantees for deep nonconvex backbones, nor does it solve distributed systems issues such as load balancing or stale concept shards. The reported fixed-bank implementation now evaluates the exact token-level CA recurrence by chunked affine scan: chunk size affects kernel tiling and parallel depth, not the mathematical causal rule. The remaining systems limitation is efficiency rather than causality: dense CA in every layer is expensive relative to highly optimized Mamba-2/RWKV kernels, so current LM prototypes allocate the dense CoM memory globally and treat layerwise or in-block CoM placements as ablations. The Concept Memory results also show that consequence-feature capacity and the learned decoder are important systems choices: too small a $d_u$ underfits the repaired consequence statistic, while the direct $d_u=d_{\rm model}$ endpoint is not automatically superior in the present residual interface. Budget maintenance in the graph-growing extension is outside the formal refinement theorem: merge, eviction, decay, and compression may discard distinctions, so their effect must be evaluated as a resource-management approximation.

The experiments are a focused sequence-modeling benchmark suite, not a full large-model scaling result. The CoM training probe in Table~\ref{tab:com-training-probe} is a same-stream training-loss comparison, not a final held-out perplexity or downstream zero-shot evaluation. It should be read as evidence for the dense CoM design direction, especially the consequence-sketch capacity result, not as evidence that CA already dominates optimized attention, Mamba, RWKV, GDN, or TTT implementations at production scale. A complete empirical evaluation should include state-equalized MQAR, language modeling at multiple parameter and token budgets, held-out same-corpus PPL, zero-shot task metrics, latency and memory curves, and ablations of chunk size, write-key shift, retention parameterization, consequence-feature dimension, fixed-bank versus dynamic concepts, and edge-free versus edge-gated variants.
